\chapter{Технологическая часть}

В данном разделе рассматривается процесс выбора инструментов для
реализации базы данных и прикладного программного обеспечения.
Обосновывается выбор системы управления базами данных и используемого
технологического стека. Описывается создание сущностей базы данных,
заданных ограничений целостности и пользовательских ролей.
Приводятся SQL-скрипты для создания таблиц и ролей в разработанной
системе, а также тестируется триггер.

\section{Выбор системы управления данными}

Для реализации системы хранения и каталогизации датасетов требуется база данных для
метаданных и отдельное хранилище для самих файлов. Для сравнения были выбраны
несколько популярных СУБД.

Сравнение выполнено по четырём простым критериям:
\begin{itemize}
  \item поддержка ролевой модели;
  \item поддержка ограничений целостности;
  \item поддержка пользовательских типов данных.
  \item наличие средств резервного копирования и восстановления.
\end{itemize}

\begin{table}[H]
\centering
\caption{Сравнение СУБД по расширенным критериям}
\label{tab:db_comparison_extended}
\renewcommand{\arraystretch}{1.2}
\begin{tabular}{|p{3cm}|c|c|c|c|c|}
\hline
\textbf{СУБД} & \textbf{Роли} & \textbf{Ограничения} & \textbf{Индексы} & \textbf{Бэкап} & \textbf{Польз. типы} \\
\hline
PostgreSQL & есть & есть & есть & есть \\
\hline
MySQL & есть & есть & есть & ограниченно \\
\hline
SQLite & нет & частично & файл & нет \\
\hline
\end{tabular}
\end{table}

\textbf{SQLite}~\cite{sqllite} --- простая встраиваемая СУБД, которая не поддерживает полноценную ролевую модель и
имеет ограниченные механизмы целостности.
Подходит для локальных приложений, но не для многопользовательской системы.

\textbf{MySQL}~\cite{mysql} является одной из наиболее распространённых СУБД и поддерживает
ролевую модель, индексы и базовые ограничения целостности. Однако её возможности по
расширяемости и созданию пользовательских типов данных значительно уступают PostgreSQL.

\noindent Таким образом, для метаданных выбран \textbf{PostgreSQL}~\cite{post} как наиболее сбалансированное решение.

Для хранения самих датасетов целесообразно использовать \textbf{объектное хранилище S3}~\cite{s3},
так как оно рассчитано на работу с большими бинарными файлами, обеспечивает высокую
надёжность хранения и простую интеграцию через стандартный протокол.

\section{Средства реализации}

Для реализации программного обеспечения системы хранения датасетов был
выбран язык программирования \textbf{Go}~\cite{go}. Данный выбор обусловлен его
строгой статической типизацией и поддержкой всех необходимых типов данных,
определённых на этапе проектирования базы данных. Язык также обеспечивает
удобные механизмы работы с многопоточностью и сетевыми операциями, что
делает его подходящим для создания серверных приложений.

Для взаимодействия между серверной частью и пользователями реализован
API на базе стандартной библиотеки Go и фреймворка
\texttt{net/http}, что предоставляет возможность
интеграции с внешними системами и клиентскими приложениями.

\section{Создание таблиц}

В данном подразделе приведены SQL-скрипты создания таблиц базы данных,
а также последующего добавления ограничений целостности.

\subsection*{Таблица Users}

\begin{lstlisting}[language=SQL,caption={Создание таблицы Users}]
CREATE TABLE users (
    id                BIGSERIAL PRIMARY KEY,
    username          VARCHAR(50)  NOT NULL,
    email             VARCHAR(255) NOT NULL,
    password          VARCHAR(255) NOT NULL,
    registration_date TIMESTAMPTZ  NOT NULL DEFAULT now(),
    country           VARCHAR(100) NOT NULL,
    is_blocked        BOOLEAN      NOT NULL DEFAULT FALSE,
    role              VARCHAR(10)  NOT NULL
);

ALTER TABLE users
    ADD CONSTRAINT uq_users_email UNIQUE (email),
    ADD CONSTRAINT chk_users_role CHECK (role IN ('guest','user','admin'));
\end{lstlisting}

\subsection*{Таблица Categories}

\begin{lstlisting}[language=SQL,caption={Создание таблицы Categories}]
CREATE TABLE categories (
    id          BIGSERIAL PRIMARY KEY,
    name        VARCHAR(50) NOT NULL,
    description TEXT
);

ALTER TABLE categories
    ADD CONSTRAINT uq_categories_name UNIQUE (name);
\end{lstlisting}

\subsection*{Таблица Datasets}

\begin{lstlisting}[language=SQL,caption={Создание таблицы Datasets}]
CREATE TABLE datasets (
    id          BIGSERIAL PRIMARY KEY,
    name        VARCHAR(100) NOT NULL,
    description TEXT,
    owner_id    BIGINT NOT NULL,
    category_id BIGINT NOT NULL,
    is_public   BOOLEAN      NOT NULL DEFAULT FALSE,
    created_at  TIMESTAMPTZ  NOT NULL DEFAULT now()
);

ALTER TABLE datasets
    ADD CONSTRAINT fk_datasets_owner FOREIGN KEY
    (owner_id) REFERENCES users(id) ON DELETE RESTRICT,
    ADD CONSTRAINT fk_datasets_category FOREIGN KEY
    (category_id) REFERENCES categories(id) ON DELETE RESTRICT;

CREATE INDEX idx_datasets_owner ON datasets(owner_id);
CREATE INDEX idx_datasets_category_public_created
    ON datasets(category_id, is_public, created_at DESC);
\end{lstlisting}

\subsection*{Таблица Dataset\_Versions}

\begin{lstlisting}[language=SQL,caption={Создание таблицы Dataset\_Versions}]
CREATE TABLE dataset_versions (
    id          BIGSERIAL PRIMARY KEY,
    number      VARCHAR(20) NOT NULL,
    upload_date TIMESTAMPTZ NOT NULL DEFAULT now(),
    filepath    TEXT        NOT NULL,
    dataset_id  BIGINT      NOT NULL,
    change_log  TEXT
);

ALTER TABLE dataset_versions
    ADD CONSTRAINT fk_versions_dataset FOREIGN KEY
    (dataset_id) REFERENCES datasets(id) ON DELETE CASCADE;

CREATE INDEX idx_versions_dataset ON dataset_versions(dataset_id, upload_date DESC);
\end{lstlisting}

\subsection*{Таблица Metadata}

\begin{lstlisting}[language=SQL,caption={Создание таблицы Metadata}]
CREATE TABLE metadata (
    id                 BIGSERIAL PRIMARY KEY,
    format             VARCHAR(50) NOT NULL,
    size               BIGINT      NOT NULL,
    tags               TEXT,
    dataset_version_id BIGINT      NOT NULL
);

ALTER TABLE metadata
    ADD CONSTRAINT chk_metadata_size CHECK (size > 0),
    ADD CONSTRAINT fk_metadata_version FOREIGN KEY
    (dataset_version_id) REFERENCES dataset_versions(id) ON DELETE CASCADE;

CREATE INDEX idx_metadata_version ON metadata(dataset_version_id);
\end{lstlisting}

\subsection*{Таблица Subscriptions}

\begin{lstlisting}[language=SQL,caption={Создание таблицы Subscriptions}]
CREATE TABLE subscriptions (
    user_id    BIGINT NOT NULL,
    dataset_id BIGINT NOT NULL,
    created_at TIMESTAMPTZ NOT NULL DEFAULT now(),
    PRIMARY KEY (user_id, dataset_id)
);

ALTER TABLE subscriptions
    ADD CONSTRAINT fk_subscriptions_user FOREIGN KEY
    (user_id) REFERENCES users(id) ON DELETE CASCADE,
    ADD CONSTRAINT fk_subscriptions_dataset FOREIGN KEY
    (dataset_id) REFERENCES datasets(id) ON DELETE CASCADE;

CREATE INDEX idx_subscriptions_dataset ON subscriptions(dataset_id);
\end{lstlisting}

\subsection*{Таблица Reviews}

\begin{lstlisting}[language=SQL,caption={Создание таблицы Reviews}]
CREATE TABLE reviews (
    id         BIGSERIAL PRIMARY KEY,
    user_id    BIGINT   NOT NULL,
    dataset_id BIGINT   NOT NULL,
    rating     SMALLINT NOT NULL,
    created_at TIMESTAMPTZ NOT NULL DEFAULT now(),
    text       TEXT
);

ALTER TABLE reviews
    ADD CONSTRAINT fk_reviews_user FOREIGN KEY
    (user_id) REFERENCES users(id) ON DELETE CASCADE,
    ADD CONSTRAINT fk_reviews_dataset FOREIGN KEY
    (dataset_id) REFERENCES datasets(id) ON DELETE CASCADE,
    ADD CONSTRAINT chk_reviews_rating CHECK (rating BETWEEN 1 AND 5);

CREATE INDEX idx_reviews_dataset ON reviews(dataset_id);
CREATE INDEX idx_reviews_user ON reviews(user_id);
\end{lstlisting}

\subsection*{Таблица Notifications}

\begin{lstlisting}[language=SQL,caption={Создание таблицы Notifications}]
CREATE TABLE notifications (
    id         BIGSERIAL PRIMARY KEY,
    user_id    BIGINT NOT NULL,
    dataset_id BIGINT NOT NULL,
    message    TEXT   NOT NULL,
    is_read    BOOLEAN NOT NULL DEFAULT FALSE,
    created_at TIMESTAMPTZ NOT NULL DEFAULT now()
);

ALTER TABLE notifications
    ADD CONSTRAINT fk_notifications_user FOREIGN KEY
    (user_id) REFERENCES users(id) ON DELETE CASCADE,
    ADD CONSTRAINT fk_notifications_dataset FOREIGN KEY
    (dataset_id) REFERENCES datasets(id) ON DELETE CASCADE;

CREATE INDEX idx_notifications_user ON notifications(user_id);
\end{lstlisting}

\section{Функция нахождения среднего рейтинга датасета}

На листинге~\ref{lst:get_avg_rating} представлена реализация функции базы данных для нахождения среднего рейтинга датасета.

\begin{lstlisting}[language=SQL,caption={Реализация функции базы данных}, label={lst:get_avg_rating}]
CREATE
OR REPLACE FUNCTION get_dataset_avg_rating(
    p_dataset_id BIGINT,
    p_return_zero_if_empty BOOLEAN DEFAULT TRUE
)
RETURNS NUMERIC(10,4)
LANGUAGE plpgsql
AS $$
DECLARE
avg_rating NUMERIC(10,4);
BEGIN
SELECT AVG(rating) ::NUMERIC(10,4)
INTO avg_rating
FROM ratings
WHERE dataset_id = p_dataset_id
  AND rating BETWEEN 1 AND 5;

IF
avg_rating IS NULL THEN
        IF p_return_zero_if_empty THEN
            RETURN 0.0;
ELSE
            RETURN NULL;
END IF;
END IF;

RETURN avg_rating;
END;
$$;
\end{lstlisting}

\section{Реализация триггера базы данных}

На листинге~\ref{lst:trig} представлена реализация триггера базы данных для автоматизации служебных действий
при создании датасета.

\begin{lstlisting}[language=SQL,caption={Реализация триггера базы данных}, label={lst:trig}]
CREATE OR REPLACE FUNCTION trg_after_insert_datasets()
RETURNS TRIGGER
LANGUAGE plpgsql
AS $$
DECLARE
v_owner_role   VARCHAR(10);
    v_owner_block  BOOLEAN;
    v_now          TIMESTAMPTZ := now();
    v_forced_msg   TEXT := NULL;
    v_placeholder  TEXT := format('pending://upload/%s', NEW.id);
BEGIN

SELECT role, is_blocked
INTO v_owner_role, v_owner_block
FROM users
WHERE id = NEW.owner_id;

IF v_owner_role IS NULL THEN
        RAISE EXCEPTION 'Owner % not found for dataset %', NEW.owner_id, NEW.id;
END IF;

    IF v_owner_block IS TRUE AND NEW.is_public IS TRUE THEN
UPDATE datasets SET is_public = FALSE WHERE id = NEW.id;
v_forced_msg := 'Dataset visibility forced to private: owner is blocked';
END IF;

INSERT INTO subscriptions(user_id, dataset_id, created_at)
VALUES (NEW.owner_id, NEW.id, v_now)
    ON CONFLICT (user_id, dataset_id) DO NOTHING;

IF NOT EXISTS (SELECT 1 FROM dataset_versions WHERE dataset_id = NEW.id) THEN
        INSERT INTO dataset_versions(number, upload_date, filepath, dataset_id, change_log)
        VALUES ('v0.1', v_now, v_placeholder, NEW.id,
                'Initial version placeholder (created by trigger)');
END IF;

INSERT INTO notifications(user_id, dataset_id, message, is_read, created_at)
VALUES (
           NEW.owner_id,
           NEW.id,
           COALESCE(
                   v_forced_msg,
                   'Dataset created. Upload the file to replace placeholder URL'
           ),
           FALSE,
           v_now
       );

RETURN NEW;
END;
$$;

DROP TRIGGER IF EXISTS trg_after_insert_datasets ON datasets;

CREATE TRIGGER trg_after_insert_datasets
    AFTER INSERT ON datasets
    FOR EACH ROW
    EXECUTE FUNCTION trg_after_insert_datasets();
\end{lstlisting}

\subsection{Тестирование триггера}

Для тестирования работы триггера были выделены классы эквивалентности, отражающие все возможные ситуации при вставке датасета.

\subsubsection*{Классы эквивалентности}

\begin{table}[H]
\centering
\caption{Классы эквивалентности для триггера \texttt{trg\_after\_insert\_datasets}}
\label{tab:trigger_eq_classes}
\begin{tabular}{|p{5cm}|p{4cm}|p{6cm}|}
\hline
\textbf{Условие} & \textbf{Класс эквивалентности} & \textbf{Ожидаемый результат} \\
\hline
Владелец существует, не заблокирован &
Нормальный ввод  &
Создаются подписка владельца на датасет и первая версия датасета.
\\
\hline
Владелец существует, но заблокирован, датасет создан публичным &
Владелец заблокирован + публичный датасет &
Поле \texttt{is\_public} принудительно устанавливается в \texttt{FALSE}, создаются подписка и версия. \\
\hline
Владелец существует, но заблокирован, датасет создан приватным &
Владелец заблокирован + приватный датасет &
Датасет остаётся приватным, создаются подписка и версия. \\
\hline
Владелец не найден в таблице \texttt{users} &
Невалидный ввод (несуществующий владелец) &
Выбрасывается исключение: \texttt{Owner <id> not found for dataset <id>}. \\
\hline
\end{tabular}
\end{table}

\section{Создание ролей}

Для обеспечения безопасности и разграничения доступа к данным
в базе данных разработана ролевая модель.
В неё входят три роли: \textit{гость}, \textit{зарегистрированный пользователь} и \textit{администратор}.
Ниже приведены SQL-скрипты, реализующие данную модель.

\subsection*{Роль Guest}

\begin{lstlisting}[language=SQL,caption={Создание роли Guest}]
CREATE ROLE guest_role WITH
    NOSUPERUSER
    NOCREATEDB
    NOCREATEROLE
    NOINHERIT
    NOREPLICATION
    NOBYPASSRLS
    CONNECTION LIMIT -1
    LOGIN
    PASSWORD 'Guestpass';

GRANT SELECT ON datasets TO guest_role;

ALTER TABLE datasets ENABLE ROW LEVEL SECURITY;

CREATE POLICY guest_datasets_select ON datasets
    FOR SELECT TO guest_role
    USING (true);
\end{lstlisting}


\subsection*{Роль Registered User}

\begin{lstlisting}[language=SQL,caption={Создание роли Registered User}]
CREATE ROLE registered_role WITH
    NOSUPERUSER
    NOCREATEDB
    NOCREATEROLE
    NOINHERIT
    NOREPLICATION
    NOBYPASSRLS
    CONNECTION LIMIT -1
    LOGIN
    PASSWORD 'Userpass';

GRANT SELECT ON datasets, versions, users TO registered_role;
GRANT INSERT ON datasets TO registered_role;
GRANT UPDATE, DELETE ON datasets TO registered_role;

ALTER TABLE datasets ENABLE ROW LEVEL SECURITY;
ALTER TABLE files ENABLE ROW LEVEL SECURITY;

CREATE POLICY user_datasets_policy ON datasets
    FOR ALL TO registered_role
    USING (owner_id = current_user::TEXT);
\end{lstlisting}


\subsection*{Роль Admin}

\begin{lstlisting}[language=SQL,caption={Создание роли Admin}]
CREATE ROLE admin_role WITH
    NOSUPERUSER
    NOCREATEDB
    NOCREATEROLE
    INHERIT
    NOREPLICATION
    NOBYPASSRLS
    CONNECTION LIMIT -1
    LOGIN
    PASSWORD 'Adminpass';

GRANT ALL PRIVILEGES ON ALL TABLES IN SCHEMA public TO admin_role;
GRANT ALL PRIVILEGES ON ALL SEQUENCES IN SCHEMA public TO admin_role;
GRANT ALL PRIVILEGES ON ALL FUNCTIONS IN SCHEMA public TO admin_role;
GRANT ALL PRIVILEGES ON ALL PROCEDURES IN SCHEMA public TO admin_role;

ALTER ROLE admin_role BYPASSRLS;
\end{lstlisting}

Таким образом, ролевая модель обеспечивает:
\begin{itemize}
    \item доступ на чтение для гостей;
    \item ограниченный доступ с возможностью модификации только собственных данных для зарегистрированных пользователей;
    \item полный контроль над всеми объектами базы данных для администраторов.
\end{itemize}

\section*{Вывод}


В технологической части был выполнен выбор средств реализации
программного обеспечения и базы данных, реализованы основные
сущности и ограничения целостности, а также разработана ролевая
модель. Приведены листинги функции и триггера базы данных,
проведено их тестирование с использованием классов эквивалентности.

